\subsection{Capacity under an equal budget (RQ2, H2)}
\label{sec:res-capacity}

Figure~\ref{fig:closed} shows closed-loop throughput as a function of the number of concurrent clients, and
Table~\ref{tab:usl} the fitted USL parameters. The monolith reached a peak of \xo{usl/browse/A/xmax} browse and
\xo{usl/checkout/A/xmax} checkout requests per second. The decomposed configurations peaked far lower: B at
\xo{xmaxratio/browse/B} (browse) and \xo{xmaxratio/checkout/B} (checkout) of the monolith's capacity, C at
\xo{xmaxratio/browse/C} and \xo{xmaxratio/checkout/C}, and D at \xo{xmaxratio/browse/D} and
\xo{xmaxratio/checkout/D}. No decomposed configuration matched the monolith at any measured concurrency level in
either workload, so the throughput crossover does not exist within the tested range and \textbf{H2 is falsified} by
its first condition. Its second condition is met only formally: the fits for B and C put $\beta$ at its lower
bound of zero, which is lower than the monolith's \xo{usl/browse/A/beta} (browse), but for the reason discussed
next rather than because decomposition removed a coherency cost.

The cgroup counters show what limits each configuration. At 64 clients the monolith's database container used
\xo{chot/browse/A/64/util}\% of its two-core limit and was throttled in \xo{chot/browse/A/64/thr}\% of its periods,
while the application container used \xo{chot/browse/A/64/apputil}\% of its four cores (checkout:
\xo{chot/checkout/A/64/util}\% and \xo{chot/checkout/A/64/apputil}\%). The monolith is therefore bound by its
database tier, and its fitted $\beta>0$ and $N^\ast$ of \xo{usl/browse/A/nstar} (browse) and \xo{usl/checkout/A/nstar}
(checkout) describe a saturated database rather than application code. In every decomposed configuration the
limit is instead the CPU quota of the edge service: \texttt{\xo{chot/browse/B/64/appname}} in~B at
\xo{chot/browse/B/64/apputil}\% of its single core and \texttt{\xo{chot/browse/C/64/appname}} in~C at
\xo{chot/browse/C/64/apputil}\% of its half core, while their databases stayed below saturation.
The edge service does the most work per request: it parses its callees' JSON responses, encodes the client's
response, and drives every outbound call, so it is the first to exhaust an equal share of the budget.

Configuration~C is the extreme case. Its USL fit has $\alpha = \xo{usl/browse/C/alpha}$ and the throughput curve is
flat from a single client onwards (\xo{x/browse/C/1} requests per second at $N=1$, \xo{x/browse/C/256} at $N=256$).
With a single client, \texttt{search} already used \xo{chot/browse/C/1/apputil}\% of its half-core limit and was
throttled in \xo{chot/browse/C/1/thr}\% of its CFS periods; additional clients only queue. In USL terms the configuration is completely contention-bound, and the
contended resource is a CPU quota. B and D have $\alpha$ between \xo{usl/checkout/D/alpha} and \xo{usl/browse/B/alpha}. D's throughput with one
client is lower than C's because of the sidecars' added latency (\xo{x/browse/D/1} against \xo{x/browse/C/1}
requests per second), but its peak is slightly higher (\xo{usl/browse/D/xmax} against \xo{usl/browse/C/xmax}). Two
measured effects explain this: the sidecars run on CPU outside the application budget, and D's edge application
container needed less CPU per request than C's (\xo{hot/browse/D/500/util}\% against \xo{hot/browse/C/500/util}\%
of the same limit at 500 requests per second).
